\begin{beginnerstatement}[Kurz erklärt: Die Architekturabbildungen lesen]

Eine Softwarearchitektur ist ein Bauplan, der Komponenten, ihre Aufgaben und
ihre Verbindungen zeigt. Eine Komponente ist dabei ein abgegrenzter technischer
Baustein. Ein Repository ist ein verwalteter Projektordner mit
Änderungshistorie; das Master-Repository ist hier die zentrale Vorlage, aus der
Projekte erzeugt werden. Der Client läuft auf der Seite der Nutzenden und enthält das sichtbare
Frontend, während der Server das Backend im Hintergrund ausführt. Eine
Datenbank speichert Daten dauerhaft und wird in passenden Profilen nur vom
Backend angesprochen. Deployment bedeutet, diese Bausteine in einer
Zielumgebung bereitzustellen und zu starten. Die Trennung von
Verantwortlichkeiten sorgt dafür, dass jeder Baustein eine klare Aufgabe hat
und möglichst eigenständig getestet oder geändert werden kann.

\end{beginnerstatement}
